\chapter{理论基础与相关技术}

\section{B/S 架构与前后端分离模式}

B/S 架构以浏览器作为统一访问入口，服务器端承担主要业务处理与数据管理职责。相较于传统 C/S 模式，B/S 架构在部署维护、跨平台访问以及多角色共享统一业务入口方面具有明显优势。对于企业供应链协同场景而言，顾客、内部管理人员与供应商往往处于不同终端和不同网络环境中，若系统需要支持灵活访问、集中维护和快速迭代，则 B/S 架构更符合实际业务需求。

前后端分离模式是在 B/S 架构基础上的进一步演进，其核心思想是将界面呈现与业务逻辑解耦，前端负责页面渲染、交互控制与状态展示，后端负责认证鉴权、业务处理、数据访问与规则校验。该模式有助于提高系统的开发效率和模块复用能力，也便于后续移动端、小程序端或第三方系统调用统一接口。结合供应链协同系统多角色、多流程、状态变更频繁的特点，前后端分离模式能够有效降低耦合度，为权限控制、消息推送和业务扩展提供良好的结构基础\citep{fielding2000architectural}。

\section{Spring Boot 框架}

Spring Boot 是基于 Spring 生态的快速开发框架，通过自动配置、约定优于配置、嵌入式容器支持以及完善的组件集成能力，简化了企业级 Web 应用的开发过程\citep{springbootdocs}。在系统实现中，开发者可围绕控制层、服务层、持久层进行分层组织，并借助统一配置文件管理安全、数据库和消息等相关参数。

对于本文设计的供应链协同管理系统而言，Spring Boot 的优势主要体现在三个方面。其一，系统业务模块较多，包括订单、库存、生产、采购、质检和消息通知等，Spring Boot 便于快速构建结构清晰的 RESTful 后端服务。其二，Spring Boot 与 Spring Security、Spring Data JPA 等组件具有良好的集成性，能够降低认证授权和数据访问模块的实现复杂度。其三，系统后期可能存在功能增量开发需求，例如报表分析、移动端接口扩展等，Spring Boot 具备较好的可维护性与扩展能力，因此适合作为本系统的后端开发框架。

\section{Spring Security 与 JWT 身份认证机制}

在多角色业务系统中，身份认证与权限控制是保障数据安全和职责边界清晰的基础。Spring Security 提供了较为完善的认证授权机制，可通过过滤器链、用户详情服务、密码编码器以及访问控制表达式等方式实现访问拦截和资源保护\citep{springsecuritydocs}。JWT 是一种基于 JSON 的令牌规范，通过对用户身份与权限声明进行签名封装，使服务器在后续请求处理中能够基于令牌完成无状态身份校验\citep{rfc7519}。

本系统采用 Spring Security 与 JWT 相结合的方式实现认证授权。用户登录成功后，服务器生成包含用户标识和角色信息的令牌，并返回给前端保存。前端在后续接口调用时将令牌附加到请求头中，后端过滤器解析令牌后完成用户身份恢复和权限上下文构建。该方式具有以下适用性：一是能够降低传统 Session 模式下的服务器状态维护压力，适合前后端分离系统；二是便于结合角色信息实现基于路径和方法级别的访问控制；三是有利于在订单审核、仓库出入库、采购审批和质检处理等高敏感操作中进行统一鉴权。

\section{Vue 3、Vite 与 Axios 前端技术}

Vue 3 是面向构建用户界面的渐进式前端框架，具备组件化开发、响应式数据驱动和良好的生态支持等特征\citep{vuedocs}。在企业供应链管理场景中，前端页面通常涉及表单录入、列表查询、状态过滤、流程操作和实时消息展示等复杂交互，Vue 3 的组件化机制能够将订单表单、库存表格、生产任务面板、采购处理页面和消息中心等视图进行模块化组织，提升前端代码复用性和维护效率。

Vite 是现代前端构建工具，依托原生 ES 模块和快速热更新机制，在开发体验和构建效率方面具有明显优势\citep{vitedocs}。Axios 则用于实现前端与后端之间的 HTTP 通信，能够统一封装请求头、响应拦截、异常处理和令牌附带逻辑。对于本系统而言，Vue 3、Vite 与 Axios 的组合不仅有利于构建响应迅速、交互明确的业务界面，也便于统一管理多角色登录态、接口异常与业务操作反馈，从而提升前端实现的规范性和一致性。

\section{WebSocket/STOMP 实时通信技术}

供应链协同系统中的多个业务节点具有明显的时效性要求。例如，销售审核结果需要及时反馈给仓库，库存不足状态需要尽快触发生产或采购处理，采购订单下发后供应商需要及时获知，质检结果也需要同步给相关岗位。若完全依赖前端轮询获取状态变化，不仅增加服务器负载，也容易带来消息延迟。WebSocket 通过在客户端与服务器之间建立持久连接，实现双向实时通信，适合处理频繁状态变化场景\citep{rfc6455}。STOMP 则在 WebSocket 之上提供了轻量级消息协议，有助于实现基于主题或队列的消息路由与订阅机制\citep{stomp12}。

本系统基于 WebSocket/STOMP 构建实时消息通知模块。后端在订单审核、库存预警、采购下发、生产任务生成和质检状态更新等关键节点发布消息，前端通过订阅相应主题获取通知并更新界面状态。该技术方案有助于减少人工催办成本，提高跨角色协同的及时性，并使系统在业务流程驱动下具备更好的状态同步能力。

\section{MySQL 与 JPA 数据持久化技术}

供应链协同管理涉及用户、角色、产品、订单、库存、生产计划、采购订单、原材料、质检记录和消息通知等大量结构化数据，系统需要在保证数据一致性的前提下完成高频读写和关联查询。MySQL 作为成熟的关系型数据库，具有较好的稳定性、事务支持能力和生态兼容性，适合承载中小型企业业务系统的核心数据\citep{mysql2024}。

JPA 提供了面向对象的数据持久化规范，使开发者能够通过实体映射方式管理数据库表结构与对象关系。基于 Spring Data JPA，可进一步简化常见的增删改查、条件查询与分页访问实现\citep{springdatajpa2024}。在本系统中，JPA 的适用性主要体现在三个方面：其一，用户、订单、库存、采购和质检等实体之间关联关系明确，适合采用实体映射建模；其二，系统需要在业务服务层组织统一事务，以确保状态流转和数据更新的一致性；其三，JPA 有助于减少重复的数据访问代码，提高后端开发效率。

\section{本章小结}

本章围绕系统设计与实现所依赖的关键技术进行了分析。B/S 架构与前后端分离模式为多角色统一访问和系统扩展提供了总体支撑；Spring Boot 为后端服务构建提供了高效框架；Spring Security 与 JWT 满足了系统对认证授权和角色控制的需求；Vue 3、Vite 与 Axios 适合构建复杂业务界面；WebSocket/STOMP 支撑了实时消息通知；MySQL 与 JPA 则为业务数据持久化提供了可靠基础。上述技术共同构成了本文企业供应链协同管理系统的技术实现基础。